\chapter{Ricerca Semantica}\label{chap:fifth}

\inlineminitoc

Il Capitolo~\ref{chap:fourth} ha descritto come i documenti normativi diventano un grafo di nodi e relazioni, i cui contenuti testuali sono rappresentati mediante embedding semantici. Questo capitolo descrive che cosa accade quando una domanda in linguaggio naturale arriva a quella base di conoscenza, e definisce con precisione le due strategie di recupero che il Capitolo~\ref{chap:sixth} mette a confronto sperimentalmente.

La distinzione fra le due strategie è deliberatamente stretta. Entrambe interrogano lo stesso grafo, indicizzano la stessa unità testuale, ricevono la stessa domanda codificata con lo stesso modello di embedding e restituiscono lo stesso numero di risultati. Differiscono per un solo elemento: se le relazioni giuridiche memorizzate nel grafo intervengano o meno nell'ordinamento dei risultati.

\section{Le due strategie e che cosa isolano}

Il sistema espone due strategie di recupero, denominate nei log e nelle tabelle dei capitoli successivi \textbf{S1} e \textbf{S2}:

\begin{itemize}
    \item \textbf{S1 --- \texttt{FullVectorSearchStrategy}}: ricerca vettoriale pura. La domanda diventa un vettore, i commi vengono ordinati per similarità coseno e i primi $k$ costituiscono il risultato. Il grafo interviene unicamente come supporto di memorizzazione: le sue relazioni non sono percorse e non influenzano l'ordinamento. Questa strategia è la \textit{baseline}.
    \item \textbf{S2 --- \texttt{LegalGraphSearchStrategy.}}: ricerca ibrida. La ricerca vettoriale produce un bacino ampio di candidati; le relazioni giuridiche del grafo producono una seconda graduatoria a partire dai primi risultati vettoriali; le due graduatorie vengono fuse. Questa è la strategia di produzione.
\end{itemize}

È importante chiarire fin d'ora una proprietà che la Sezione~\ref{sec:aritmetica_fusione} motiverà quantitativamente: \textbf{nella configurazione adottata, S2 agisce sul risultato finale come un \textit{re-ranker} della ricerca vettoriale}. Il ramo giuridico può raggiungere ulteriori commi, ma il peso assegnato al solo segnale strutturale non è sufficiente a farli entrare nella Top-10. Il contributo del grafo consiste quindi soprattutto nel promuovere candidati già presenti nel bacino vettoriale. Questa non è una limitazione accidentale, ma una scelta di progetto motivata nella Sezione~\ref{sec:aritmetica_fusione}.

Nessuna delle due strategie riceve informazioni sul tipo di domanda, sui documenti attesi o sulle etichette del benchmark. Il modulo di recupero riceve esclusivamente la stringa della domanda; le etichette esistono solo dentro il modulo di valutazione e non attraversano mai il confine. Questa separazione evita che il recupero utilizzi direttamente le risposte attese e rende possibile un confronto controllato fra le strategie.

\section{S1: la ricerca vettoriale pura}
\label{sec:ricerca_vettoriale_pura}

La strategia di base realizza il percorso descritto in forma generale nella Sezione~\ref{sec:creazione_embedding}. La domanda viene codificata con il medesimo modello impiegato per indicizzare i commi --- condizione necessaria, perché due modelli diversi producono spazi vettoriali non confrontabili --- e il vettore risultante viene sottoposto a Neo4j tramite l'indice vettoriale del nodo \texttt{:Paragraph}.

\begin{lstlisting}[language=Cypher, caption={Interrogazione dell'indice vettoriale per la ricerca di base}, label={lst:cypher_vector}]
CALL db.index.vector.queryNodes($index_name, $top_k, $query_vector)
YIELD node, score
MATCH (act:LegalAct)-[:HAS_TITLE|HAS_CHAPTER|HAS_ARTICLE*1..4]->
      (art:Article)-[:HAS_PARAGRAPH]->(node)
RETURN node.custom_id AS paragraph_custom_id,
       act.title AS atto, art.number AS num_articolo,
       node.number AS num_comma, node.text AS testo, score
ORDER BY score DESC
\end{lstlisting}

L'indice restituisce i vicini approssimati del vettore di query; la risalita gerarchica successiva serve unicamente ad arricchire ciascun risultato con la sua provenienza documentale, affinché il risultato restituito all'utente sia citabile. Non modifica l'ordine.

Il costo computazionale è quello dell'\textit{Approximate Nearest Neighbor Search}: sui 6.396 commi del corpus la latenza mediana misurata è di circa 16~ms per interrogazione, esclusa la codifica della domanda.

Questa strategia incarna un'ipotesi precisa sul dominio: che la pertinenza di un comma rispetto a una domanda sia catturabile dalla vicinanza semantica fra il testo della domanda e il testo del comma. È un'ipotesi ragionevole e, come mostrerà il Capitolo~\ref{chap:sixth}, spesso sufficiente. Il suo limite emerge quando la risposta completa richiede una disposizione che il comma pertinente \textit{richiama} senza riprodurne il contenuto. In quel caso il comma richiamato può essere lessicalmente e semanticamente distante dalla domanda: un encoder può ancora recuperarlo attraverso segnali presenti nel testo, ma non dispone di una rappresentazione esplicita del rinvio che ne giustifica la pertinenza.

\section{S2: la ricerca arricchita dalla struttura}

La strategia ibrida è organizzata in cinque step spiegati brevemente con lo schema seguente, e poi descritti in dettaglio nelle sottosezioni successive.

\begin{center}
\texttt{Top-200 vettoriale $\rightarrow$ Top-5 semi $\rightarrow$ relazioni giuridiche $\rightarrow$ RRF pesato $\rightarrow$ testa riservata}
\end{center}

\subsection{Step 1: il bacino vettoriale}

Il primo step è identico a S1, ma con un cutoff molto più ampio: vengono recuperati i \textbf{primi 200} commi per similarità coseno. Questa lista, denominata \texttt{vector\_pool}, costituisce il bacino dei candidati sostenuti dalla ricerca vettoriale e la base sulla quale il ramo giuridico esercita il proprio effetto di riordinamento.

La scelta di un bacino ampio ha una motivazione precisa. Il compito del grafo non è trovare commi che la ricerca vettoriale non vede affatto, ma \textit{promuovere} commi che la ricerca vettoriale vede male: un comma che il coseno colloca in posizione 30 o 120, troppo in basso per essere restituito, ma che una relazione giuridica raggiunge a partire da un risultato di testa. Un bacino di 200 elementi rende visibili al meccanismo di fusione proprio quei casi.

\subsection{Step 2: i semi del ramo giuridico}

I \textbf{primi 5} commi del bacino vettoriale vengono designati come \textit{semi} (\textit{seeds}) dell'espansione a grafo. Sono gli unici nodi da cui le relazioni verranno percorse.

Il numero è piccolo, e lo è per una ragione sperimentale. L'implementazione originaria impiegava 20 semi, e il risultato era un peggioramento sistematico: una citazione che parte dal ventesimo comma migliore non è un indizio, è rumore. Il ramo giuridico ha senso solo se i punti da cui parte sono essi stessi affidabili, perché il suo compito è raggiungere ciò che è \textit{collegato a una risposta corretta}, non ciò che è collegato a un candidato qualsiasi.

\subsection{Step 3: il ramo giuridico}

Da ciascun seme viene percorso, con un'unica interrogazione Cypher, l'insieme delle relazioni giuridiche descritte nella Sezione~\ref{sec:estrazione_relazioni}: \texttt{CITES}, \texttt{DEROGATES}, \texttt{INTERPRETS}, \texttt{REPEALS} e \texttt{REPLACES}. Le relazioni sono percorse in \textbf{entrambe le direzioni}: se un comma cita una disposizione, quella disposizione è pertinente al comma, ma è vero anche il contrario --- una norma successiva che deroga alla norma trovata è altrettanto necessaria per rispondere correttamente.

Ogni comma raggiunto riceve un punteggio pari al coseno del seme da cui proviene, moltiplicato per un coefficiente che dipende dal tipo di relazione:

\begin{table}[H]
    \centering
    \renewcommand{\arraystretch}{1.2}
    \begin{tabular}{|l|c|l|}
        \hline
        \textbf{Relazione} & \textbf{Moltiplicatore} & \textbf{Natura del legame} \\
        \hline
        \texttt{REPEALS}, \texttt{REPLACES} e \texttt{DEROGATES} & 1,00 & modifica l'efficacia della norma \\
        \texttt{INTERPRETS} & 0,90 & vincola l'interpretazione \\
        \texttt{CITES} & 0,70 & rinvio generico \\
        \hline
    \end{tabular}
    \caption{Moltiplicatori per tipo di relazione giuridica nel ramo a grafo.}
    \label{tab:moltiplicatori_relazione}
\end{table}

La gerarchia riflette la forza giuridica del legame: una disposizione che abroga o deroga alla norma trovata è \textit{necessaria} per una risposta corretta, mentre una citazione generica è soltanto \textit{pertinente}.

La lista risultante viene deduplicata e troncata a \texttt{top\_k} $\times$ 2 righe.


\subsection{Step 4: la fusione RRF pesata}
\label{sec:aritmetica_fusione}

Le due graduatorie --- quella vettoriale e quella giuridica --- vengono combinate mediante una variante pesata della \textit{Reciprocal Rank Fusion}. Per ciascun comma \(d\):

\begin{equation}
\mathrm{score}(d)
=
\sum_{i \in \{\text{vett},\text{leg}\}}
\frac{w_i}{k_{\mathrm{RRF}}+\mathrm{rank}_i(d)}
\label{eq:rrf}
\end{equation}

con \(k_{\mathrm{RRF}}=60\), \(w_{\mathrm{vett}}=4\) e \(w_{\mathrm{leg}}=3\). Un comma presente in una sola graduatoria riceve un solo contributo; se compare in entrambe, i due contributi vengono sommati.

La RRF opera sui \textit{ranghi} anziché sui punteggi originari delle due graduatorie. Questa proprietà è particolarmente utile nel caso considerato, poiché evita di confrontare direttamente grandezze di natura diversa, come la similarità coseno della ricerca vettoriale e il punteggio prodotto dal ramo giuridico.

I valori dei pesi non sono stati fissati arbitrariamente. A parità di \(k_{\mathrm{RRF}}\), infatti, l'ordinamento prodotto dalla fusione dipende soltanto dal rapporto

$$
\lambda =
\frac{w_{\mathrm{leg}}}{w_{\mathrm{vett}}},
$$

poiché moltiplicare entrambi i pesi per una stessa costante riscala tutti i punteggi RRF senza modificarne l'ordine. La scelta può quindi essere trattata come la selezione di un unico iperparametro.

Il rapporto è stato sottoposto a un'analisi di sensibilità sul benchmark di sviluppo, mantenendo invariata la restante configurazione della pipeline. L'analisi non ha individuato un optimum puntuale, ma una regione di stabilità nella quale variazioni del rapporto non producono modifiche sostanziali alla graduatoria finale. Il valore

$$
\lambda = 0{,}75,
$$

equivalente ai pesi \(4:3\), è stato scelto come punto interno a tale regione e successivamente mantenuto fisso nella valutazione finale. Non va quindi interpretato come un valore calibrato con precisione, ma come una configurazione rappresentativa di un intervallo stabile.

La scelta soddisfa inoltre un vincolo strutturale della strategia. Con \(w_{\mathrm{vett}}=4\), il miglior punteggio ottenibile da un comma presente esclusivamente nella graduatoria giuridica è

$$
\frac{w_{\mathrm{leg}}}{61},
$$

mentre un comma collocato al decimo posto della graduatoria vettoriale riceve

$$
\frac{4}{70}.
$$

Affinché un candidato esclusivamente giuridico possa superare quest'ultimo deve quindi valere

$$
\frac{w_{\mathrm{leg}}}{61}
\geq
\frac{4}{70},
$$

da cui

$$
w_{\mathrm{leg}}
\geq
4\frac{61}{70}
\simeq
3{,}49.
$$

Il valore adottato, \(w_{\mathrm{leg}}=3\), rimane al di sotto di questa soglia. Il ramo giuridico può quindi promuovere un comma che possiede anche evidenza vettoriale, ma il solo segnale strutturale non è sufficientemente forte da far entrare nella Top-10 un candidato assente dalla graduatoria vettoriale. Questa proprietà mantiene S2, nella configurazione adottata, nel ruolo di \textit{re-ranker} della ricerca vettoriale, anziché trasformarla in una seconda strategia di recupero indipendente.

La Tabella~\ref{tab:aritmetica_rrf} mostra il comportamento concreto della formula.

\begin{table}[H]
\centering
\renewcommand{\arraystretch}{1.2}
\begin{tabular}{|l|r|l|}
\hline
\textbf{Condizione del comma} & \textbf{Punteggio RRF} & \textbf{Esito} \\
\hline
solo vettoriale, rango 1 & 0,06557 & resta in testa \\
solo vettoriale, rango 10 & 0,05714 & ultimo ammesso \\
solo vettoriale, rango 30 & 0,04444 & fuori dal Top-10 \\
vettoriale rango 30 + rango giuridico 1 & \textbf{0,09362} & \textbf{promosso in testa} \\
solo giuridico, rango 1 & 0,04918 & resta fuori dal Top-10 \\
\hline
\end{tabular}
\caption{Comportamento della fusione RRF pesata con \(k_{\mathrm{RRF}}=60\) e pesi \(4:3\).}
\label{tab:aritmetica_rrf}
\end{table}

L'esempio evidenzia il comportamento desiderato della fusione: il segnale giuridico può modificare in modo marcato la posizione di un comma già sostenuto dalla ricerca vettoriale, mentre da solo non è sufficiente a sostituire i candidati presenti nella parte utile della graduatoria semantica.

\subsection{Step 5: la testa riservata alla vettoriale}
\label{sec:testa_riservata}

La fusione, per come è costruita, premia i commi che compaiono in entrambe le liste. Questo è l'obiettivo della strategia, ma comporta un rischio specifico: il comma che risponde \textit{direttamente} alla domanda --- quello che il coseno aveva collocato al primo posto --- può essere scavalcato da un comma collegato, che il grafo ha promosso. Per una risposta completa questo è irrilevante, perché entrambi sono nel risultato; ma per un sistema a valle che legge i risultati nell'ordine in cui arrivano, non lo è affatto.

Il meccanismo che previene questo effetto è la \textbf{testa riservata}: i primi $\mathrm{round}(k \times 0{,}2)$ ranghi del risultato finale vengono riportati in puro ordine di coseno. Sono 2 posizioni in un Top-10 e 1 in un Top-5. Le righe sono quelle prodotte dalla fusione, con il loro punteggio e la loro provenienza; soltanto l'ordine viene ripristinato.

Una conseguenza del meccanismo è che, nelle posizioni appartenenti alla testa riservata, i punteggi RRF non sono necessariamente disposti in ordine decrescente. Un comma promosso dalla fusione e collocato immediatamente al di sotto della testa può infatti possedere un punteggio RRF superiore a quello di uno dei risultati mantenuti nelle prime posizioni. Non si tratta di un'incoerenza nell'ordinamento: in quelle posizioni prevale deliberatamente l'ordine vettoriale sul punteggio prodotto dalla fusione.

L'effetto di questo componente sull'ordinamento viene quantificato nel Capitolo~\ref{chap:sixth}. La testa riservata introduce infatti uno scambio esplicito: limita la libertà della fusione nei primi ranghi per preservare in cima i candidati con la maggiore evidenza semantica diretta.


\subsection{Un ramo rimosso: l'adiacenza fra commi}

Accanto al ramo giuridico ha operato per un certo periodo un secondo ramo, che percorreva la relazione \texttt{NEXT\_PARAGRAPH} per raggiungere i commi contigui a quelli trovati. L'intuizione è forte: se la risposta sta nel comma 3 di un articolo, è probabile che il comma 4 la completi.

Le prove preliminari condotte sui pesi storicamente impiegati hanno mostrato un deterioramento sistematico delle prestazioni al crescere del contributo di questo ramo. L'adiacenza è stata pertanto esclusa dalla configurazione sottoposta alla valutazione finale, che mantiene soltanto le relazioni giuridiche esplicite.

Nel complesso, le scelte descritte in questo capitolo definiscono le due strategie di recupero e isolano il contributo delle relazioni giuridiche al riordinamento dei risultati. Il Capitolo~\ref{chap:sixth} presenta i risultati dell'esperimento e mette a confronto quantitativamente le prestazioni delle due strategie.
